\documentclass[11pt,twocolumn]{article}

% ------------------------------------------------------------
% Page layout
% ------------------------------------------------------------
\usepackage[letterpaper,margin=0.75in]{geometry}
\usepackage{microtype}
\usepackage{parskip}
\setlength{\columnsep}{0.25in}

% ------------------------------------------------------------
% Fonts: use LuaLaTeX or XeLaTeX
% ------------------------------------------------------------
\usepackage{fontspec}
\usepackage{unicode-math}

\setmainfont{DejaVu Serif}
\setsansfont{DejaVu Sans}
\setmonofont{DejaVu Sans Mono}
\setmathfont{Latin Modern Math}

% ------------------------------------------------------------
% Core math
% ------------------------------------------------------------
\usepackage{amsmath}
\usepackage{amssymb}
\usepackage{mathtools}
\usepackage{bm}
\usepackage{tensor}
\usepackage{cancel}
\usepackage{mathrsfs}

% ------------------------------------------------------------
% Tables
% ------------------------------------------------------------
\usepackage{booktabs}
\usepackage{array}
\usepackage{tabularx}
\usepackage{multirow}
\usepackage{siunitx}

\sisetup{
  detect-all,
  per-mode=symbol,
  exponent-product=\cdot
}

% ------------------------------------------------------------
% Figures, SVG, PGF/TikZ plots
% ------------------------------------------------------------
\usepackage{graphicx}
\usepackage{svg}
\usepackage{tikz}
\usepackage{pgfplots}
\usepackage{stfloats}

\pgfplotsset{compat=1.18}

% Optional PGF libraries
\usetikzlibrary{
  arrows.meta,
  calc,
  positioning,
  decorations.pathmorphing,
  matrix
}

% ------------------------------------------------------------
% Captions and references
% ------------------------------------------------------------
\usepackage{caption}
\usepackage{subcaption}
\usepackage[hidelinks]{hyperref}
\usepackage[nameinlink,noabbrev]{cleveref}

% ------------------------------------------------------------
% Common math macros
% ------------------------------------------------------------

% Sets and spaces
\newcommand{\R}{\mathbb{R}}
\newcommand{\C}{\mathbb{C}}
\newcommand{\N}{\mathbb{N}}
\newcommand{\Z}{\mathbb{Z}}
\newcommand{\Q}{\mathbb{Q}}

\newcommand{\Ltwo}{L^2}
\newcommand{\Lp}[1]{L^{#1}}
\newcommand{\Hsob}[1]{H^{#1}}
\newcommand{\Cinfty}{C^\infty}

% Scalars, vectors, matrices, tensors
\newcommand{\scal}[1]{#1}
\newcommand{\vct}[1]{\bm{#1}}
\newcommand{\mat}[1]{\bm{#1}}
\newcommand{\ten}[1]{\bm{\mathsf{#1}}}

% Unit vectors
\newcommand{\ihat}{\hat{\vct{i}}}
\newcommand{\jhat}{\hat{\vct{j}}}
\newcommand{\khat}{\hat{\vct{k}}}

% Transpose, inverse, trace, determinant
\newcommand{\T}{^{\mathsf{T}}}
\newcommand{\inv}{^{-1}}
\newcommand{\tr}{\operatorname{tr}}
\newcommand{\diag}{\operatorname{diag}}
\newcommand{\rank}{\operatorname{rank}}
\newcommand{\nullspace}{\operatorname{null}}
\newcommand{\range}{\operatorname{range}}

% Norms, inner products
\newcommand{\norm}[1]{\left\lVert #1 \right\rVert}
\newcommand{\abs}[1]{\left\lvert #1 \right\rvert}
\newcommand{\inner}[2]{\left\langle #1, #2 \right\rangle}

% Differential operators
\newcommand{\dd}{\,\mathrm{d}}
\newcommand{\pd}{\partial}
\newcommand{\grad}{\nabla}
\newcommand{\divg}{\nabla \cdot}
\newcommand{\curl}{\nabla \times}
\newcommand{\rot}{\nabla \times}
\newcommand{\laplace}{\Delta}

% Derivatives
\newcommand{\dv}[2]{\frac{\mathrm{d} #1}{\mathrm{d} #2}}
\newcommand{\dvtwo}[2]{\frac{\mathrm{d}^2 #1}{\mathrm{d} #2^2}}
\newcommand{\pdv}[2]{\frac{\partial #1}{\partial #2}}
\newcommand{\pdvtwo}[2]{\frac{\partial^2 #1}{\partial #2^2}}

% Matrix exponential
\newcommand{\expm}{\operatorname{expm}}

% Complex / imaginary numbers
\newcommand{\ii}{\mathrm{i}}
\newcommand{\Real}{\operatorname{Re}}
\newcommand{\Imag}{\operatorname{Im}}

% Fourier transform
\newcommand{\FT}{\mathcal{F}}
\newcommand{\IFT}{\mathcal{F}^{-1}}
\newcommand{\fourier}[1]{\widehat{#1}}

% Probability and statistics
\newcommand{\E}{\mathbb{E}}
\newcommand{\Var}{\operatorname{Var}}
\newcommand{\Cov}{\operatorname{Cov}}
\newcommand{\Prob}{\mathbb{P}}
\newcommand{\Normal}{\mathcal{N}}

% Quaternions
\newcommand{\quat}[1]{\mathbf{#1}}
\newcommand{\qconj}[1]{#1^{*}}
\newcommand{\qinv}[1]{#1^{-1}}
\newcommand{\qmul}{\otimes}

% Stochastic calculus
\newcommand{\Wiener}{W_t}
\newcommand{\ito}{\mathrm{d}}

% Numerical methods
\newcommand{\bigO}{\mathcal{O}}

% ------------------------------------------------------------
% Document metadata
% ------------------------------------------------------------
\title{Two-Column Mathematical Notes}
\author{Mikhail Grushinskiy}
\date{\today}

\begin{document}

\maketitle

\begin{abstract}
A compact two-column template for mathematical writing with support for
vector algebra, matrix methods, calculus, quaternions, tensors, Fourier
analysis, PDEs, SDEs, numerical methods, tables, PGF charts, and SVG graphics.
\end{abstract}

\section{Vector and Matrix Algebra}

Vectors are written in bold:
\[
\vct{x} =
\begin{bmatrix}
x_1 \\ x_2 \\ x_3
\end{bmatrix},
\qquad
\mat{A} =
\begin{bmatrix}
a_{11} & a_{12} \\
a_{21} & a_{22}
\end{bmatrix}.
\]

A linear system is
\[
\mat{A}\vct{x} = \vct{b},
\qquad
\vct{x} = \mat{A}\inv \vct{b}.
\]

The matrix exponential is
\[
\vct{x}(t) = \exp(\mat{A}t)\vct{x}_0
           = \expm(\mat{A}t)\vct{x}_0.
\]

\section{Calculus}

A gradient, divergence, curl, and Laplacian are written as
\[
\grad f,
\qquad
\divg \vct{u},
\qquad
\curl \vct{u},
\qquad
\laplace f.
\]

Surface and volume integrals:
\[
\iint_{\partial \Omega} \vct{F} \cdot \vct{n}\,\dd S
=
\iiint_{\Omega} \divg \vct{F}\,\dd V.
\]

\section{Quaternions}

A quaternion can be written as
\[
\quat{q}
=
q_0 + q_1\ihat + q_2\jhat + q_3\khat.
\]

Quaternion rotation may be written as
\[
\vct{v}' = \quat{q} \qmul \vct{v} \qmul \qconj{\quat{q}}.
\]

\section{Statistics}

A Gaussian random vector:
\[
\vct{x} \sim \Normal(\vct{\mu}, \mat{\Sigma}).
\]

Mean and covariance:
\[
\E[\vct{x}] = \vct{\mu},
\qquad
\Cov[\vct{x}] = \E[(\vct{x}-\vct{\mu})(\vct{x}-\vct{\mu})\T].
\]

\section{Fourier Analysis}

The Fourier transform is written as
\[
\fourier{f}(\omega)
=
\FT\{f(t)\}
=
\int_{-\infty}^{\infty}
f(t)e^{-\ii \omega t}\,\dd t.
\]

The inverse transform is
\[
f(t)
=
\IFT\{\fourier{f}(\omega)\}
=
\frac{1}{2\pi}
\int_{-\infty}^{\infty}
\fourier{f}(\omega)e^{\ii \omega t}\,\dd \omega.
\]

\section{PDEs and SDEs}

A diffusion PDE:
\[
\pdv{u}{t}
=
\alpha \laplace u.
\]

An Itô SDE:
\[
\ito X_t
=
a(X_t,t)\,\ito t
+
b(X_t,t)\,\ito W_t.
\]

\section{Infinite Series and Limits}

A Taylor expansion:
\[
f(x)
=
\sum_{n=0}^{\infty}
\frac{f^{(n)}(a)}{n!}(x-a)^n.
\]

A limit:
\[
\lim_{h \to 0}
\frac{f(x+h)-f(x)}{h}
=
\dv{f}{x}.
\]

\section{Tensors}

A second-order tensor may be written as
\[
\ten{T}
=
T_{ij}\,\vct{e}_i \otimes \vct{e}_j.
\]

Using the \texttt{tensor} package:
\[
\tensor{T}{^i_j}
=
g^{ik}T_{kj}.
\]

\section{Numerical Methods}

Forward Euler:
\[
y_{n+1}
=
y_n + h f(t_n,y_n)
+
\bigO(h^2).
\]

Runge--Kutta notation:
\[
k_1 = f(t_n,y_n),
\qquad
k_2 = f\left(t_n+\frac{h}{2}, y_n+\frac{h}{2}k_1\right).
\]

\section{Tables}

\begin{table}[h]
\centering
\caption{Example numerical results.}
\begin{tabular}{lcc}
\toprule
Method & Error & Runtime \\
\midrule
Euler & \num{1.2e-2} & \SI{0.3}{ms} \\
RK4   & \num{2.1e-5} & \SI{1.4}{ms} \\
\bottomrule
\end{tabular}
\end{table}

\section{PGF Chart}

\begin{figure}[h]
\centering
\begin{tikzpicture}
\begin{axis}[
    width=\columnwidth,
    height=0.65\columnwidth,
    xlabel={$x$},
    ylabel={$f(x)$},
    grid=major
]
\addplot[
    domain=0:6.28,
    samples=100
]
{sin(deg(x))};
\end{axis}
\end{tikzpicture}
\caption{A PGFPlots sine curve.}
\end{figure}

\section{SVG Graphics}

Place an SVG file in the same folder, for example \texttt{device.svg}, then use:

\begin{figure}[h]
\centering
% Requires: latexmk -lualatex -shell-escape main.tex
% Requires Inkscape installed.
\includesvg[width=0.8\columnwidth]{device}
\caption{Example SVG graphic.}
\end{figure}

\section{Wide Two-Column Figure}

Use \texttt{figure*} for a figure spanning both columns.

\begin{figure*}[t]
\centering
\begin{tikzpicture}
\begin{axis}[
    width=0.85\textwidth,
    height=0.28\textwidth,
    xlabel={$t$},
    ylabel={$x(t)$},
    grid=major
]
\addplot[
    domain=0:10,
    samples=200
]
{exp(-0.2*x)*sin(deg(2*x))};
\end{axis}
\end{tikzpicture}
\caption{A wide two-column PGF chart.}
\end{figure*}

\end{document}
